\documentclass{article}

\PassOptionsToPackage{numbers, sort, compress}{natbib}
\usepackage[main,final]{neurips_2025}

\usepackage[T2]{fontenc}
\usepackage[utf8]{inputenc}
\usepackage[english,russian]{babel}
\usepackage{hyperref}
\usepackage{url}
\usepackage{booktabs}
\usepackage{nicefrac}
\usepackage{microtype}
\usepackage{xcolor}

%%%

\usepackage{subcaption}
\usepackage{graphicx}
\usepackage{multirow}
\usepackage{amsmath,amssymb,amsfonts}
\usepackage{amsthm}
\usepackage{mathrsfs}
\usepackage{xcolor}
\usepackage{textcomp}
\usepackage{manyfoot}
\usepackage{algorithm}
\usepackage{algorithmicx}
\usepackage{algpseudocode}
\usepackage{listings}

\newtheorem{theorem}{Theorem} % continuous numbers
%%\newtheorem{theorem}{Theorem}[section] % sectionwise numbers
%% optional argument [theorem] produces theorem numbering sequence instead of independent numbers for Proposition
\newtheorem{proposition}[theorem]{Proposition}% 
\newtheorem{lemma}{Lemma}% 
%%\newtheorem{proposition}{Proposition} % to get separate numbers for theorem and proposition etc.

\newtheorem{example}{Example}
\newtheorem{remark}{Remark}

\newtheorem{definition}{Definition}
\newtheorem{assumption}{Assumption}

%%%


\title{Инвариантный индекс текста на основе собственной размерности при машинном переводе}


% The \author macro works with any number of authors. There are two commands
% used to separate the names and addresses of multiple authors: \And and \AND.
%
% Using \And between authors leaves it to LaTeX to determine where to break the
% lines. Using \AND forces a line break at that point. So, if LaTeX puts 3 of 4
% authors names on the first line, and the last on the second line, try using
% \AND instead of \And before the third author name.


\author{
  Marina Kuchina\\
  MIPT\\
  Moscow, Russia\\
  \texttt{kuchina.ma@phystech.edu}\\
  \And
  Ivan Papay\\
  MIPT\\
  Moscow, Russia\\
  \texttt{papai.id@phystech.edu}\\
  \AND
  Andrii Hraboviy\\
  MIPT\\
  Moscow, Russia\\
  \texttt{email@phystech.edu}\\
  \And
  Alexander Kildyakov\\
  MIPT\\
  Moscow, Russia\\
  \texttt{email@phystech.edu}\\
}


\begin{document}


\maketitle

\begin{abstract}
  В данной работе решается задача поиска близких документов в мультиязычных коллекциях. Для её решения исследуется инвариантность собственной размерности текста при переводе. Также проверяется возможность использования собственной размерности для построения поискового индекса документа. Исследование основано на вычислении собственной размерности для мультиязычных текстов одной тематики и анализе полученных значений. В работе собственная размерность впервые рассматривается как представление текста, устойчивое к переводу. В результате подтверждена инвариантность собственной размерности при переводе и показана возможность построения на её основе поискового индекса, позволяющего находить тематически схожие документы независимо от языка.
\end{abstract}

\textbf{Keywords:} Intrinsic dimension, Persistent homology dimension, Cross-lingual retrieval, Translation invariance, Natural language processing, Text embeddings

\section{Введение}\label{sec:intro}

В наши дни растут объёмы информации в мультиязычных цифровых коллекциях, и задача эффективного поиска информации, не ограниченной языковыми барьерами, становится критически важной [Bakhteev, 2019]. Одной из фундаментальных проблем в этой области является поиск семантически близких документов на разных языках. Многие из существующих подходов опираются на параллельные корпуса текстов на разных языках и сложные модели машинного перевода. Альтернативный подход заключается в поиске характеристик текста, инвариантных относительно перевода и отражающих его смысл. В данной работе рассматривается внутренняя размерность текста – характеристика сложности его внутренней структуры. Задача заключается в проверке инвариантности внутренней размерности текста относительно перевода и построении на её основе индекса текста, который можно было бы использовать для поиска похожих по смыслу документов в мультиязычных коллекциях.

В литературе представлено несколько подходов к решению задачи межъязыкового поиска документов. Более ранние подходы основываются на построении общих эмбеддинговых пространств при помощи кросс-язычных словарей или параллельных корпусов текстов [Vulić & Moens, 2015]. Другие методы опираются на количественные меры семантической близости текстов [Wäschle, 2015]. Современный подход заключается в использовании многоязычных моделей-энкодеров, например, XLM-RoBERTa [Conneau et al., 2020] или RoBERTa [Liu et al., 2019], проецирующих мультиязычные тексты в единое векторное пространство и позволяющие сравнивать полученные представления.

Однако существующие методы обладают рядом ограничений. Так, например, подход  [Vulić & Moens, 2015], основанный на построении общих эмбеддинговых пространств, строит модель для двух языков и, как следствие, плохо масштабируется на случай мультиязычных коллекций. В случае использования энкодерных моделей ряд современных работ, посвящённых исследованию геометрии их представлений, показывает, что эмбеддинги текстов содержат чувствительные к языку компоненты [Chang et al., 2022], а также зависят от контекста [Ethayarajh, 2019]. Это доказывает, что сравнение эмбедингов в них напрямую может дать нестабильные результаты, так как смысловая информация искажается специфичной для языка и контекстуальной информацией.

В данной работе исследуется внутренняя размерность текста на инвариантность относительно перевода. Для оценки внутренней размерности текста используется метод, основанный на теории персистентной гомологии [Barannikov, 1994], а именно оценка персистентной размерности (PHD) [Adams et al., 2020; Schweinhart, 2021]. Множество эмбеддингов фрагмента текста рассматривается как облако точек. Предполагается, что точки лежат на некотором многообразии, размерность которого количественно определяет сложность текста. Метод PHD оценивает размерность многообразия через анализ минимального остовного дерева (MST), построенного на этих точках. Сумма длин рёбер MST, возведённых в степень $\alpha$, растёт с числом точек по закону $n^{(d - \alpha)/d}$, где d — искомая размерность [Steele, 1988; Birdal et al., 2021]. Анализ роста суммы с ростом числа точек и подбор параметра $\alpha$ позволяет восстановить искомую размерность d. В современных работах [Tulchinskii et al., 2023] уже показано, что PHD инвариантна относительно перефразирования. В нашей работе проверяется инвариантность PHD относительно перевода. На основе этой гипотезы строится индекс документа как последовательность значений PHD для его фрагментов. Полученная последовательность считается идентификатором текста для поиска в многоязычной коллекции документов.

Для проверки гипотезы об инвариантности PHD относительно перевода использован датасет Wiki-40B [Guo et al., 2020]. Для статей на одну тему, рассматриваемых в качестве приближения перевода, были получены эмбеддинги и вычислено значение PHD. Анализ полученных значений подтвердил высокую корреляцию между размерностями для статей на разных языках, что свидетельствует об инвариантности внутренней размерности PHD при переводе. При разбиении текста на фрагменты и вычислении последовательностей PHD для полученных фрагментов был построен поисковый индекс, который позволил находить тематически схожие документы независимо от языка. Результаты свидетельствуют о перспективности использования индексов текстов, основанных на PHD, для поиска текстов в мультиязычных коллекциях.

\section{Related Work}\label{sec:rw}

\textbf{Topic \#1.}
TODO

\textbf{Topic \#2.}
TODO

\section{Preliminaries}\label{sec:prelim}

\subsection{General notation}

In this section, we introduce the general notation used in the rest of the paper and the basic assumptions. 

\subsection{Assumptions} 

TODO

\section{Method}\label{sec:method}

\section{Experiments}\label{sec:exp}

To verify the theoretical estimates obtained, we conducted a detailed empirical study...

\section{Discussion}\label{sec:disc}

TODO

\section{Conclusion}\label{sec:concl}

TODO


%%%%%%%%%%%%%%%%%%%%%%%%%%%%%%%%%%%%%%%%%%%%%%%%%%%%%%%%%%%%

\bibliographystyle{unsrtnat}
\bibliography{references}

%%%%%%%%%%%%%%%%%%%%%%%%%%%%%%%%%%%%%%%%%%%%%%%%%%%%%%%%%%%%

\newpage
\appendix
\section{Appendix / supplemental material}\label{app}

\subsection{Additional experiments / Proofs of Theorems}\label{app:exp}

TODO

\end{document}